\begin{abstract}
The Lightning Network routes payments over chains of payment channels using
hash time-locked contracts (HTLCs). We present one \tamarin{} theory covering
the channel layer (unboundedly many updates, revocation, punishment) and
multi-hop routing over routes of any length, against a Dolev-Yao adversary who
can also compromise long-term keys. We prove that only cheating or
compromised parties are punished, that settlements are authenticated by
the receiver's invoice, and that preimages, which our model never sends on the
network, stay secret. Under three liveness assumptions, imposed as restrictions
that remove every run violating them, an uncompromised
intermediary never has its outgoing HTLC redeemed while its incoming one is
refunded, and payment atomicity holds in runs without key compromise. The
assumption that intermediaries claim in time is necessary: without it,
\tamarin{} finds a run without compromise in which an intermediary loses
funds. We give the wormhole attack as a \tamarin-checked trace: colluding
intermediaries bypass an honest one, which is refunded but earns no fee.
Conversely, we prove for routes of any length that a hop of a settled payment
whose hash is offered once can be bypassed only if some long-term key is
compromised. On a fixed route, per-hop locks in the style of point time-locked
contracts (PTLCs) prevent the wormhole around an honest intermediary when the
sender is honest, whereas plain HTLCs admit it. All results are machine-checked
with \tamarin~1.10 and reproduced with one command.
\end{abstract}

\begin{IEEEkeywords}
Lightning Network, payment channels, HTLC, wormhole attack, symbolic
verification, \tamarin
\end{IEEEkeywords}
